% Generated draft: requires mathematical and source review.
\documentclass[11pt]{article}
\usepackage[margin=25mm]{geometry}
\usepackage{fontspec}
\setmainfont{Latin Modern Roman}
\usepackage{amsmath,amssymb,mathtools,unicode-math}
\setmathfont{Latin Modern Math}
\usepackage{enumitem,needspace,xurl,hyperref}
\setlength{\parindent}{0pt}
\setlength{\parskip}{5pt}
\setlength{\emergencystretch}{3em}
\newcommand{\N}{\mathbb{N}}
\newcommand{\Z}{\mathbb{Z}}
\newcommand{\Q}{\mathbb{Q}}
\newcommand{\R}{\mathbb{R}}
\newcommand{\C}{\mathbb{C}}
\newcommand{\F}{\mathbb{F}}
\newcommand{\A}{\mathbb{A}}
\newcommand{\PP}{\mathbb{P}}
\newcommand{\E}{\mathbb{E}}
\newcommand{\eps}{\varepsilon}
\newcommand{\dd}{\,\mathrm{d}}
\begin{document}
\section*{UnsolvedMath ID: 30003519}
\textbf{Source problem code:} catalogue entry
\subsection{Definitions and mathematical statement}
A monomorphic population has a single resident heritable trait; its spatial population profile need not be constant. A monomorphic equilibrium is a stationary population profile with that trait. A jump process of monomorphic equilibria describes transitions between such equilibria. The supplied record refers to such a process but does not specify its state space, transition rates, or underlying spatial population model.

\textbf{Requested mathematical statement.} For the spatial adaptive-dynamics model intended by the source, derive a canonical equation describing the evolution of the resident trait in the limit of asymptotically small mutation steps, starting from the jump process of monomorphic equilibria.

For notation only, let \(X^{\varepsilon}\) denote the resident-trait jump process when mutation steps have size controlled by \(\varepsilon>0\). The requested derivation would identify an appropriate time rescaling \(a_{\varepsilon}>0\), a limiting trait trajectory \(x\), and an evolution law, schematically
\[
X^{\varepsilon}(a_{\varepsilon}\,\cdot)\longrightarrow x(\cdot),
\qquad
\frac{\mathrm d x}{\mathrm d t}=b(x(t)),
\qquad \varepsilon\downarrow0.
\]
This display describes the intended type of canonical-equation limit, not a fully specified convergence assertion. The trait space, initial conditions, convergence mode, time interval, mutation law, scaling, and vector field \(b\) are absent from the supplied record. The spatial domain, boundary conditions, ecological dynamics, and hypotheses ensuring existence and suitable selection of resident equilibria are also unspecified. Consequently, no precise model-specific theorem or universal quantification over spatial models can be extracted from this record alone.
\subsection{Short English statement}
Derive a small-mutation-step canonical equation from the spatial jump process of monomorphic equilibria intended by the source. The supplied record labels this problem open, but does not provide the model or hypotheses needed for a precise convergence theorem. This is a statement-only draft requiring human mathematical review.
\subsection{Sources}
\begin{itemize}
\item Supplied local UnsolvedMath registry record: permanent ID \(30003519\), title \emph{Canonical Equations for Spatial Adaptive Dynamics}, source problem code \texttt{catalogue entry}, in the Oberwolfach Reports Open Problems collection. No individual proposer is supplied; the recorded proposed year is 2017.
\item Supplied source citation: \emph{Reaction Networks and Population Dynamics (2018)}, Oberwolfach Reports, pp.~1747--1804, DOI \texttt{10.4171/owr/2017/28}. Original report link: \url{https://doi.org/10.4171/owr/2017/28}. The supplied background describes these reports as documenting mathematical workshops and their open-problem discussions.
\item Supplied literature reference: Champagnat, Jabin, and Meleard, described as a structured model: \url{https://doi.org/10.1016/j.matpur.2013.10.003}.
\item Supplied literature reference: Champagnat and Hass, described as canonical-equation convergence: \url{https://arxiv.org/abs/2303.09936}.
\item Literature assessment checked at \texttt{2026-08-21}: this is supplied dated curation, not independently verified. It labels the problem open and states that rigorous canonical-equation derivations exist for several well-mixed and specially structured models, while the general spatial limit requested here remains unsettled. It further states that the cited work obtains a jump-process limit under restrictive assumptions rather than the desired spatial canonical equation. The supplied background calls this dated literature triage derived from OWR curation and requests independent verification of time-sensitive status and cited claims.
\end{itemize}
\end{document}
